\subsection{Plattform \enquote{Resteria}}\label{sec:konzept}

Der Name \enquote{Resteria} verbindet \enquote{Reste} mit der Endung \enquote{-ia}, die oft einen Ort oder eine Gemeinschaft bezeichnet; Resteria ist damit als Ort gedacht, nicht nur als bloßes Suchwerkzeug. Neben der gesellschaftlichen Bedeutung, die die Einleitung herausstellt, verfolgt Resteria ein wirtschaftliches Ziel über ein Freemium-Modell: Die Kernfunktionen bleiben kostenlos; ein kostenpflichtiges Zusatzpaket erweitert die Menüplanung um eine Vorratsübersicht mit Haltbarkeitsdaten und blendet Werbung aus. Kostenpflichtig sind dabei bewusst nur Komfortfunktionen der individuellen Planung; Reste-Matching, Rettungs-Feed und Challenges bleiben kostenlos, weil hinter einer Bezahlschranke zu wenige Beiträge zusammenkämen, um die Community zu tragen. Hinzu kommen Kooperationen mit regionalen Food-Sharing-Initiativen und Lebensmittelhändler:innen, etwa gesponserte Challenges oder Rabattaktionen für gerettete Zutaten.

Die einzelnen Funktionen folgen einer Unterscheidung aus der Verhaltensforschung zu Lebensmittelverschwendung. \textcite[S. 105]{vittuariHowReduceConsumer2023} trennen \enquote{Drivers}, also psychologische und soziale Einflussfaktoren, die Verhalten erklären, von \enquote{Levers}, den konkret nutzbaren Hebeln, mit denen sich Verhalten gezielt verändern lässt. Resteria setzt an der zweiten Kategorie an, denn jede Funktion übersetzt einen bekannten Hebel in ein konkretes Feature, statt allgemein auf Nachhaltigkeit zu setzen.

Die Kernfunktion ist das Reste-Matching. Nutzer:innen geben die Zutaten ein, die sie noch verwerten möchten, und die Plattform schlägt passende Rezepte vor. Ein anfänglicher Rezeptbestand aus Rezeptkooperationen und lizenzierten Sammlungen ist damit Voraussetzung für den Beta-Test in~\autoref{sec:phasenplanung}; Nutzerbeiträge erweitern ihn erst danach. Diese Funktion allein unterscheidet Resteria noch nicht von Restegourmet oder der \ac{BMLEH}-App, wie~\autoref{sec:wettbewerbsanalyse} gezeigt hat.

Um aus einmaliger Nutzung wiederkehrendes Verhalten zu machen, bemessen sich Rettungspunkte nach Menge und Art der verwerteten Reste-Zutat. Sie kumulieren zu einem Reste-Level; ab bestimmten Schwellen erhalten Nutzer:innen Titel wie \enquote{Resteheld:in}, die auf dem Profil sichtbar sind. Diese Mechanik ist eine Form von Gamification nach dem allgemeinen Nudge-Prinzip, also einem Gestaltungsimpuls, der Verhalten lenkt, ohne Wahlmöglichkeiten einzuschränken \parencite[S. 2]{barkerWhatNudgeTechniques2021}, ergänzt um eigene Fortschrittsanzeigen. Der Punktezähler liefert direktes Feedback zum Fortschritt, den das Reste-Level in gröberen Stufen zusammenfasst.

Ergänzend plant Resteria eine wiederkehrende Erinnerungsbenachrichtigung, etwa \enquote{Deine Reste warten}, wenn längere Zeit keine Zutat eingetragen wurde. \textcite[S. 10]{barkerWhatNudgeTechniques2021} identifizieren genau solche \enquote{Reminders} als einen der wenigen Nudge-Typen, die sich in methodisch hochwertigen Studien als wirksam erwiesen haben. \textcite[S. 14]{nivedhithaHowGreenSocial2024} zeigen zudem, dass eine grüne Selbstwahrnehmung das Verhalten nicht unmittelbar verändert, sondern erst über die gemeinsame Handlungsabsicht in der Gruppe; der Titel wirkt deshalb vor allem über die Challenge-Rangliste aus~\autoref{sec:community_und_feedback}, die Nutzer:innen sichtbar vergleicht.

Dass Gamification als Hebel grundsätzlich wirkt, zeigt eine Feldstudie von \textcite[S. 9]{somaFoodWasteReduction2020}: Vielspieler:innen einer zwölfwöchigen Intervention reduzierten ihren Lebensmittelabfall um 51 Prozent, Wenigspieler:innen nur um 39 Prozent; nach der Kampagne war der Abstand zwischen beiden Gruppen signifikant. Die Studie stützt damit die Entscheidung für Gamification, belegt aber nicht die Wirkung der konkreten Punkte- und Level-Mechanik von Resteria, weil sie eine andere Intervention getestet hat. Deren Ausgestaltung bleibt eine eigene Konzeptleistung dieser Arbeit.

Nutzer:innen können eigene Reste-Kreationen per Foto oder kurzem Video hochladen, allerdings mit Pflichtangabe der verwendeten Reste-Zutat und optional einer kurzen Zubereitungsbeschreibung in Textform. Ein Bewertungs- und Kommentarsystem erlaubt es, auf fremde Beiträge zu reagieren und Zubereitungstipps auszutauschen. Personalisierte Empfehlungen schlagen Rezepte vor, die zu bereits genutzten Zutaten oder bevorzugten Küchenstilen passen. Eine vierte Funktion verknüpft Menüplanung mit dem Zero-Waste-Fokus: Ein Wochenplan verteilt die im Reste-Matching vorgeschlagenen Rezepte über die Woche und ist über die Startseite erreichbar; das zugehörige Reste-Journal protokolliert die verwerteten Zutaten im Profilbereich, wo sie als Rettungspunkte und als Einsparangabe sichtbar werden.

Gegenüber den untersuchten Plattformen ist neu, wie Resteria den Zero-Waste-Fokus in die von der Aufgabenstellung genannten Zwecke Rezept-Teilen und gegenseitige Inspiration hineinträgt und ihn zugleich über die einmalige Nutzung hinaus verankert: Die verpflichtende Zutaten-Kennzeichnung macht das Rezept-Teilen selbst zur Verwertungshandlung, die monatlich wiederkehrende Zero-Waste-Woche ersetzt Gruppen und Foren als Anlass für gegenseitige Inspiration, und die Kopplung von Menüplanung und Reste-Journal an die Rettungspunkte bindet die Reste-Verwertung in die Wochenplanung ein. Rezept-Upload, Bewertungs- und Kommentarsystem sowie personalisierte Empfehlungen übernimmt Resteria dagegen bewusst als Formate aus etablierten Rezept-Communitys, weil sie sich dort bereits bewährt haben.~\autoref{tab:funktionsuebersicht} fasst die Funktionen mit ihrer Herleitung zusammen.

\begin{table}[H]
\caption{Funktionsübersicht von Resteria}\label{tab:funktionsuebersicht}
{\footnotesize
\begin{tabularx}{\textwidth}{@{}l X X@{}}
\toprule
\textbf{Funktion} & \textbf{Kurzbeschreibung} & \textbf{Herleitung/Bezug} \\
\midrule
Reste-Matching & Zutaten-Eingabe führt zu passenden Rezeptvorschlägen & Kernfunktion, schließt die Marktlücke aus~\autoref{sec:wettbewerbsanalyse} \\
Rettungspunkte \& Reste-Level & Punkte nach Menge/Art verwerteter Zutaten, kumulieren zu Leveln/Titeln & geplante Erinnerungsbenachrichtigung \parencite[S. 10]{barkerWhatNudgeTechniques2021}, Feedback/Fortschritt als eigene Gestaltung, Selbstwahrnehmung über gemeinsame Absicht \parencite[S. 14]{nivedhithaHowGreenSocial2024}, Gamification-Effekt \parencite[S. 9]{somaFoodWasteReduction2020} \\
\shortstack[l]{Rezept-Upload\\(Foto/Video, Text optional)} & Eigene Reste-Kreationen teilen & Aufgabenstellung: Rezept-Teilen \\
Bewertungs-/Kommentarsystem & Feedback und Zubereitungstipps austauschen & Aufgabenstellung: Kochtipp-Austausch \\
Personalisierte Empfehlungen & Rezeptvorschläge nach Nutzungsverhalten & Aufgabenstellung: Inspiration \\
Menüplanung \& Reste-Journal & Wochenplan verknüpft mit Protokoll verwerteter Zutaten & Zero-Waste-Kernmechanismus \\
\bottomrule
\end{tabularx}
}
\quelle{Eigene Darstellung.}
\end{table}

Das User Interface (UI) folgt drei Hauptbereichen: einer Startseite mit Reste-Eingabefeld und Rezeptvorschlägen, dem Rettungs-Feed aus~\autoref{sec:community_und_feedback} sowie einem Profilbereich mit Rettungspunkten, Reste-Level und eigenen Beiträgen. Über ein Lesezeichen-Icon lassen sich Rezeptvorschläge als \enquote{Gespeicherte Rezepte} merken. Die User Experience (UX) setzt auf kurze Wege vom Rest zum Rezept, weshalb Vorschläge nach Zubereitungszeit und fehlenden Zutaten sortiert sind. Farbwelt, Typografie und wiederkehrende Bedienelemente hält der Style Tile in Anhang B fest,~\autoref{fig:styletile}. Darauf aufbauend liegen für die drei Bereiche vier Mockups in Anhang C vor: die Startseite in~\autoref{fig:mockup_start}, dieselbe mit ausgeklappten gespeicherten Rezepten in~\autoref{fig:mockup_start_bookmarks}, der Rettungs-Feed in~\autoref{fig:mockup_feed} und das Profil in~\autoref{fig:mockup_profil}.

Resteria verzichtet auf einen Event-Kalender, weil ein Kalender für Kochkurse oder lokale Treffen den Fokus von der individuellen Resteverwertung auf Veranstaltungsorganisation verschieben würde und eigene Anforderungen an lokale Vernetzung und Terminpflege stellt – anders als der rein individuelle Wochenplan. Vernetzung verlagert Resteria stattdessen in den Rettungs-Feed und die Challenges.
